\chapter{Conclusiones}

El desarrollo del presente trabajo práctico permitió abordar de forma integral la cadena de diseño, acondicionamiento y
mitigación de perturbaciones en sistemas de instrumentación electrónica industrial. La articulación entre el
relevamiento de maquinaria real en los laboratorios de mecánica y la experimentación sistemática en el banco de ensayos
eléctricos consolidó los criterios ingenieriles requeridos para la adquisición confiable de señales débiles.

A partir de los experimentos y mediciones realizadas se desprenden las siguientes conclusiones fundamentales, dando
respuesta a los interrogantes planteados:

\section{Interferencias de Línea de 50 Hz}
La presencia de perturbaciones acopladas desde la red eléctrica comercial ($\qty{220}{\volt}$ a $\qty{50}{\hertz}$)
es omnipresente e inevitable en cualquier entorno residencial, universitario o fabril. Todo conductor eléctrico o
traza de circuito impreso actúa involuntariamente como antena ante los intensos campos eléctricos y magnéticos
generados por las canalizaciones de potencia. La alta impedancia de entrada característica del instrumental y de
las etapas de amplificación ($\qty{1}{\mega\ohm}$ o superior) provoca que corrientes parásitas de orden nanoamperio
originen caídas de tensión de cientos de milivoltios o voltios (como se verificó en el Experimento 3 con valores
de hasta $\qty{3.2}{\volt}$ pico a pico). En consecuencia, el sistema de adquisición no puede asumir un entorno
electromagnético ideal, resultando imperioso proyectar apantallamientos conductores, filtros pasabajos o de muesca
(\emph{notch}) y etapas diferenciales con alta relación de rechazo de modo común (CMRR $> 90\si{\deci\bel}$).

\section{Identificación de Otras Fuentes de Interferencia}
En las mediciones sobre el banco de celda de carga se identificaron claramente dos tipologías de perturbación:
\begin{enumerate}
    \item \emph{Interferencia electromagnética transitoria de banda ancha (EMI radiada)}: Producida intencionalmente
    por las chispas de conmutación en el colector de delgas del motor universal (Dremel). Este fenómeno genera ráfagas
    de impulsos de elevadísima frecuencia (decenas de megahercios) que se propagan tanto por radiación directa a través
    del espacio como por conducción a lo largo de los conductores de red.
    \item \emph{Ruido de fondo térmico y fluctuaciones ambientales}: Observado en el piso de ruido aleatorio de
    $\qty{8.57}{\milli\volt}_{\text{rms}}$ en la salida directa de la celda, generado por la agitación térmica de los
    portadores en las resistencias del puente de galgas y ruido de baja frecuencia ($1/f$).
\end{enumerate}

\section{Propuestas de Mejora en el Sistema Completo}
Para optimizar el desempeño y garantizar la compatibilidad electromagnética (EMC) del sistema frente a entornos
industriales agresivos, se proponen las siguientes mejoras:
\begin{itemize}
    \item \emph{Selección de Componentes}: Reemplazar el amplificador por dispositivos integrados con inmunidad
    reforzada contra interferencias de radiofrecuencia (\emph{EMI Hardened Op-Amps}, tales como el INA828 o AD8421),
    los cuales incorporan filtros pasabajos internos en sus etapas de entrada que atenúan las componentes de alta
    frecuencia antes de que alcancen los transistores activos. Incorporar además perlas de ferrita
    (\emph{ferrite beads}) y supresores de tensión transitoria (diodos TVS bidireccionales) en cada línea de señal.
    \item \emph{Diseño de Placa de Circuito Impreso (PCB)}: Implementar placas de al menos cuatro capas, destinando un
    plano interno continuo e ininterrumpido exclusivamente a plano de masa (retorno de referencia). Trazar las líneas de
    entrada diferencial ($V_{\text{in}}^+$, $V_{\text{in}}^-$) en estricta simetría geométrica (par diferencial con
    longitudes y anchos idénticos) para maximizar la cancelación de ruido en modo común. Separar físicamente la masa
    analógica de bajo nivel (AGND), la masa digital (DGND) y la masa de chasis/potencia (PGND), unificándolas en un
    único punto común en estrella (\emph{star grounding}). Disponer capacitores cerámicos de desacoplo de muy baja
    inductancia parásita ($\qty{100}{\nano\farad}$ MLCC) inmediatamente adyacentes a los pines de alimentación de
    cada circuito integrado.
    \item \emph{Armado de Cables y Blindajes}: Utilizar cables de par trenzado apantallado (STP - \emph{Shielded Twisted
    Pair}). El trenzado regular minimiza el área neta del bucle conductor, anulando por cancelación geométrica las
    fuerzas electromotrices inducidas por campos magnéticos alternos, mientras que la malla metálica continua confina
    las líneas de campo electrostático. A bajas frecuencias, la malla debe conectarse a tierra en un único extremo para
    prevenir la formación de bucles de masa (\emph{ground loops}); para inmunidad en radiofrecuencia, se emplean
    conectores metálicos con continuidad de apantallamiento en $360^\circ$ hacia el gabinete conductor.
\end{itemize}

\section{Normativa y Guías Técnicas de Mitigación}
El análisis experimental desarrollado se fundamenta y complementa con la literatura técnica especializada:
\begin{itemize}
    \item La nota de aplicación de Texas Instruments \emph{SLOA034: EMI Demodulation in Operational Amplifiers}
    \cite{ti_sloa034} describe en profundidad el mecanismo físico causante del corrimiento de continua
    ($\qty{+800}{\milli\volt}$) observado durante el ensayo con el motor Dremel. Establece que las junturas PN
    de entrada actúan como rectificadores no lineales frente a picos de RF, prescribiendo la inclusión de filtros
    $RC$ simétricos de modo común y diferencial para mantener la señal dentro de la región lineal.
    \item La guía de Tacuna Systems \emph{Reducing Noise in Very Sensitive Load Cell Applications} \cite{tacuna_noise}
    recomienda el empleo de conexiones de 6 hilos con líneas remotas de sensado (\emph{sense}) para compensar caídas de
    tensión en cables largos, y el uso de fuentes de excitación ratiométricas donde la tensión del ADC acompaña las
    derivas de la alimentación del puente.
    \item Tratados de referencia como el de Henry Ott \cite{ott_emc} detallan la teoría de compatibilidad
    electromagnética, señalando que la eficacia de un blindaje reside críticamente en una baja impedancia de conexión a
    tierra, tal como se corroboró al comparar la malla derivada a masa frente a la condición de malla flotante.
\end{itemize}

\section{Criterio Decisivo en la Selección del Sistema Embebido}
Al seleccionar un sistema embebido para instrumentar magnitudes mecánicas (fuerza, deformación, presión y torque),
el factor de mayor peso técnico radica en la \emph{Resolución Efectiva (ENOB), la Estabilidad y el Ruido del Subsistema
Analógico (ADC)}, antes que en la frecuencia de reloj del procesador o la cantidad de núcleos de cómputo.

Dado que las magnitudes mecánicas varían a velocidades bajas o moderadas (requiriendo frecuencias de muestreo
típicamente inferiores a $\qty{100}{\text{S/s}}$), una arquitectura con conversor sigma-delta de
$\qtyrange{16}{24}{\text{bits}}$, referencia de tensión de bajísima deriva térmica ($< 10\text{ ppm/}^\circ\text{C}$) y
rechazo intrínseco de $\qty{50}{\hertz} / \qty{60}{\hertz}$ resulta decisivamente superior a un microcontrolador de alta
velocidad dotado de un ADC SAR de 12 bits ruidoso cuyos últimos 2 o 3 bits quedan inutilizados por fluctuaciones
espurias.
